\section{Banc d'essai sur un marché Totem synthétique}\label{sec-benchmark}

\subsection{Conception}
Le marché synthétique est construit à partir des contributeurs, pour que chaque mécanisme de la \cref{sec-anatomy}
apparaisse comme dans Totem et non comme une forme ajoutée à une courbe.

\paragraph{Catégories.} Quatre catégories, onze échéances Totem d'un mois à cinq ans et six périodes de dividendes (la
fin de 2026, puis 2027 à 2031) ; date de valorisation au 14 octobre 2026.
\begin{center}\input{generated/\lang/tab_categories}\end{center}

\paragraph{Contributeurs.} Les panels options et dividendes sont tirés indépendamment. Chaque contributeur a des vues
cohérentes (\cref{ass-consistent}) : des dividendes avec un facteur persistant plus un bruit propre à chaque dividende,
dont la dispersion croît avec le temps, et une courbe de repo lisse (niveau plus inclinaison). Les contributeurs options
valorisent calls et puts à cinq strikes avec leur propre forward, une nappe de volatilité avec smile, un bruit de
volatilité et un bruit de prix. Les panels se réduisent avec la maturité et sont emboîtés.

\paragraph{Consensus et chaîne de traitement.} Le consensus est la moyenne après un filtre médiane/MAD ; dispersions et
nombres de contributeurs sont publiés comme le ferait Totem. La chaîne du contrôleur s'exécute ensuite sans
modification : forwards par parité, allocation sur le profil de détachement FO, bootstrap, covariance, détection.

\paragraph{Marquages FO.} Dividendes vrais avec un bruit de $1\,\%$ et repo vrai avec une erreur de niveau de $3$\,pb (un
écart-type).

\paragraph{Scénarios.} Chaque sous-jacent de test est simulé sous chaque scénario avec la même graine : les différences
entre scénarios sont appariées.
\begin{center}\small
\begin{tabular}{p{4.0cm}p{7.4cm}p{3.2cm}}
\toprule
Scénario & Injection & Amplitude\\
\midrule
Décalage du spot & les contributeurs options valorisent avec un autre spot & $\pm0{,}03\,\%$ à $\pm0{,}15\,\%$\\
Erreur de période & une période (2027--2029) erronée dans le consensus dividendes & $\pm0{,}15$ à $\pm0{,}5$ pour $100$ de spot\\
Désaccord de date & $35$ à $70\,\%$ du panel options placent le dividende le plus proche d'une échéance de l'autre côté & un dividende\\
Profil périmé & profil d'allocation décalé pour qu'un dividende franchisse une échéance & $3$ à $20$ jours au-delà\\
Forward aberrant & calls de consensus décalés à une échéance & $3$ à $6$ écarts-types du bruit de portage\\
Sous-panel long biaisé & les contributeurs restant après la plus forte sortie sont biaisés sur les dividendes suivants & $\pm8\,\%$ à $\pm20\,\%$\\
Erreur de marquage FO & repo FO décalé parallèlement, Totem intact & $\pm15$ à $\pm50$\,pb\\
Biais brut/net & toutes les périodes Totem multipliées & $\pm8\,\%$ à $\pm18\,\%$\\
\bottomrule
\end{tabular}
\end{center}

\paragraph{Détecteurs.} (1) Écart naïf : plus grand $|\Delta b|$ sur les piliers. (2) Forme naïve : plus grand écart du
repo implicite Totem à l'interpolation de ses deux voisins. (3) B-spline plus mélange gaussien sur la courbe Totem :
coefficients d'une B-spline cubique ajustée par moindres carrés (six coefficients, nœuds en $\ln T$), regroupés par un
mélange gaussien par catégorie dont le nombre de composantes est choisi par BIC ; le score est l'opposé de la
log-vraisemblance. (4) La même chose sur l'écart Totem--FO. (5) Le détecteur structurel de la \cref{sec-detection}.

\paragraph{Protocole.} \vNTrain{} sous-jacents sains par catégorie calibrent chaque détecteur à \vFPR{} de fausses
alertes (un quantile du score pour les références, $\kappa$ et l'échelle de niveau pour le détecteur structurel).
\vNTestPerCat{} nouveaux sous-jacents par catégorie sont ensuite simulés sous chacun des \vNScen{} scénarios, soit
\vNTest{} sous-jacents de test. Le plan étant apparié, le rapport signal sur bruit d'une injection se calcule exactement,
$\lambda=\|PL^{-1}(y^{\text{scénario}}-y^{\text{sain}})\|$ : c'est la grandeur du \cref{cor-power}.

\subsection{Résultats}
\begin{table}[ht]
\centering
\caption{Part des sous-jacents signalés, chaque détecteur étant calibré à \vFPR{} de fausses alertes sur l'historique
sain. La ligne « sain » est le taux de fausses alertes réalisé sur de nouveaux sous-jacents. Pour l'erreur de marquage FO
et le biais brut/net, une alerte n'est pas le résultat souhaité (voir le texte). Colonnes : écart naïf, forme naïve,
B-spline plus mélange gaussien sur la courbe et sur l'écart Totem--FO, détecteur structurel.}\label{tab-detection}
\input{generated/\lang/tab_detection}
\end{table}

\begin{figure}[ht]
\centering\includegraphics[width=\linewidth]{figures/\lang/detection.pdf}
\caption{Part des sous-jacents signalés par scénario et par détecteur (données du \cref{tab-detection}).}
\label{fig-detection}
\end{figure}

\begin{figure}[ht]
\centering\includegraphics[width=\linewidth]{figures/\lang/power.pdf}
\caption{Taux de détection en fonction du rapport signal sur bruit $\lambda$ de l'anomalie injectée, sur l'ensemble des six
anomalies côté Totem ; $n$ est le nombre de sous-jacents par classe. La courbe grise est la puissance théorique d'un
gabarit connu au seuil familial typique (\cref{cor-power}).}
\label{fig-power}
\end{figure}

\begin{table}[ht]
\centering
\caption{Aire sous la courbe ROC par rapport aux sous-jacents sains (qualité du classement, indépendante du seuil).}\label{tab-auc}
\input{generated/\lang/tab_auc}
\end{table}

\begin{table}[ht]
\centering
\caption{Détecteur structurel : part des sous-jacents signalés par scénario et par catégorie.}\label{tab-bycat}
\input{generated/\lang/tab_bycat}
\end{table}

\begin{figure}[ht]
\centering\includegraphics[width=\linewidth]{figures/\lang/attribution.pdf}
\caption{Mécanisme cité en premier par le détecteur structurel, par scénario. Niveau, marche, pic et saison sont les
familles de gabarits ; « rugosité » signifie que le test global s'est déclenché sans gabarit ; « pas d'alerte » signifie
que le sous-jacent n'a pas été signalé.}
\label{fig-attribution}
\end{figure}

\begin{table}[ht]
\centering
\caption{Détecteur structurel : détection, attribution (la première explication cite le mécanisme injecté et sa
position) et alertes d'erreur de marquage FO.}\label{tab-attribution}
\input{generated/\lang/tab_attribution}
\end{table}

\begin{table}[ht]
\centering
\caption{Taux de fausses alertes réalisé sur de nouveaux sous-jacents sains, par catégorie (cible \vFPR), et $\kappa$
calibré.}\label{tab-fpr}
\input{generated/\lang/tab_fpr}
\end{table}

\subsection{Lecture des résultats}
\begin{itemize}
\item \textbf{La détection moyenne ne départage pas les meilleures méthodes.} Sur les six anomalies côté Totem, le
détecteur structurel signale \vMeanStruct{} des sous-jacents, la B-spline plus mélange gaussien sur l'écart
\vMeanGmmDiff, sur la courbe \vMeanGmmCurve, le contrôle de forme naïf \vMeanShape{} et l'écart naïf \vMeanNaive.
\item \textbf{Les forwards aberrants} sont le cas où l'espace de portage et le modèle de bruit rapportent le plus :
\vFwdStruct{} détectés, contre \vFwdNaive{} pour l'écart naïf, qui ne peut pas voir une valeur aberrante de trois à six
écarts-types à une échéance longue sous le bruit du long terme, et \vFwdShape{} pour le contrôle de forme naïf.
\item \textbf{Les décalages de spot} (\vSpotStruct) et \textbf{les profils de détachement périmés} (\vSeasonStruct) sont
faciles pour toutes les méthodes ; le détecteur structurel ajoute l'explication (\cref{fig-attribution}).
\item \textbf{Les erreurs de période} (\vBucketStruct{} contre \vBucketNaive{} pour l'écart naïf) sont limitées par leur
taille : une erreur de $0{,}15$ à $0{,}5$ pour $100$ de spot est faible devant la dispersion des dividendes lointains, et la
\cref{fig-power} montre que les cas manqués ont un faible rapport signal sur bruit.
\item \textbf{Les désaccords de date de détachement} sont détectés dans \vTimingStruct{} des cas ; le contrôle de forme
naïf (\vTimingShape) est compétitif car ces pics sont grands dans l'espace de repo.
\item \textbf{Un sous-panel long biaisé} est presque non identifiable pour toutes les méthodes (\vPanelStruct{} pour le
détecteur structurel, \vPanelNaive{} pour l'écart naïf) : le biais s'accumule de façon lisse après la sortie de
contributeurs et se projette surtout sur une pente de repo. C'est le \cref{thm-nonid} à l'œuvre, pas un problème de
réglage.
\item \textbf{Les véritables erreurs de marquage FO} restent distinctes des courbes Totem aberrantes : \vFoStruct{}
déclenchent une anomalie (le taux sain) et \vFoMismark{} sont signalées comme erreurs de marquage FO possibles, contre
\vFoClean{} des sous-jacents sains. L'écart naïf et le mélange gaussien sur l'écart déclenchent une alerte de courbe
aberrante sur \vFoNaive{} et \vFoGmmDiff{} d'entre elles. Le test d'erreur de marquage FO trouve \vFoMismarkEuStock{} des
erreurs sur les actions européennes et \vFoMismarkIndex{} sur l'indice, mais seulement \vFoMismarkHtb{} sur les titres
difficiles à emprunter : les vues de repo des contributeurs y diffèrent d'environ $60$\,pb, si bien qu'un écart FO de
$15$ à $50$\,pb reste dans l'incertitude du consensus lui-même.
\item \textbf{Le biais brut/net} est signalé dans \vGrossStruct{} des cas, surtout sous forme de marches aux dates de
détachement pour les payeurs discrets ; pour le flux de dividendes d'un indice, il ressemble à une pente de repo et
ressort via $\beta$ (\cref{prop-gross}).
\item \textbf{Puissance.} Entre $4$ et $6$ de rapport signal sur bruit, le détecteur structurel trouve \vPowStructE{} des
anomalies, et au-delà de $6$ de \vPowStructF{} à \vPowStructH, conformément à la théorie du \cref{cor-power}. En dessous de
$3$, l'écart naïf et le mélange gaussien sur l'écart détectent davantage (\vPowGmmDiffC{} contre \vPowStructC{} entre $2$
et $3$) : ils réagissent aussi à la part d'une anomalie qui ressemble à un écart de repo lisse, que le détecteur
structurel laisse volontairement au test d'erreur de marquage FO. C'est cette même propriété qui leur fait déclencher des
alertes sur les véritables erreurs de marquage FO.
\item \textbf{Attribution.} Quand le détecteur structurel signale une anomalie Totem, sa première explication cite le
mécanisme injecté et sa position dans \vAttr{} des cas.
\end{itemize}

\subsection{Limites des éléments présentés}
\begin{caveat}
Le marché synthétique reproduit les mécanismes par construction ; il ne montre pas que les données Totem réelles ne
contiennent que ceux-là. Les données réelles comporteront aussi des contributeurs figés, des opérations sur titres et des
dividendes exceptionnels, de l'exercice américain dans les cotations d'actions et des erreurs de saisie. L'échelle $\kappa$
absorbe un niveau de bruit mal spécifié, mais pas une structure mal spécifiée, et les amplitudes, bien que plausibles, ne
sont calibrées sur aucun historique particulier. Le banc d'essai montre que la méthode fonctionne quand les mécanismes sont
présents et qu'elle est calibrée quand ils sont absents ; la validation sur douze à vingt-quatre mois d'instantanés réels
avec des incidents étiquetés est l'étape suivante.
\end{caveat}
